\chapter{Clocking, Timing Analysis, and Multiple Clock Domains}
\label{ch:clocks}
\chaptermark{Clocking and timing analysis}

A synchronous circuit works correctly only if every signal manages to travel
from one flip-flop to the next within a single clock cycle. In
Chapter~\ref{ch:sync} we derived the condition for the maximum frequency. Now
we will see how software checks this condition, and what to do when a circuit
has more than one clock signal or when a signal arrives from outside.

\section{The Timing Condition in an FPGA}

For a pair of flip-flops with combinational logic between them, the
condition

\[
T \;\ge\; t_{cq} + t_{\text{logic}} + t_{\text{wire}} + t_{su}
\]

must hold, where $t_{cq}$ is the delay of the first flip-flop's output after
the edge, $t_{\text{logic}}$ is the logic delay (LUT), $t_{\text{wire}}$ is the
interconnect delay, and $t_{su}$ is the setup time of the second flip-flop. In
an FPGA, wire delay makes up a significant part of the total time, so it
cannot be estimated until the circuit has been placed on the chip.

That is why, after place and route, nextpnr performs \term{static timing
analysis}: it goes through all the paths between flip-flops, adds up the
delays, and compares the result with the clock period. The difference between
the period and the length of the longest path is the \term{slack}. It is
positive if the circuit fits in the period and negative if it does not. The
longest path is called the \term{critical path}, and the highest frequency
that its delay allows is called $f_{max}$.

\begin{figure}[H]
\centering
\begin{tikzpicture}[font=\sffamily\small, x=1cm, y=1cm]
  \draw[->] (-0.2,0) -- (11,0) node[right] {time};
  \foreach \x/\t in {0/{edge 1}, 8/{edge 2}} {
    \draw[FpgaBlue, thick] (\x,-0.3) -- (\x,2.8); \node at (\x,3.1) {\t};
  }
  \fill[FpgaLight, draw=FpgaBlue] (0,0.2) rectangle (1.2,0.8) node[midway] {$t_{cq}$};
  \fill[FpgaLight!70!white, draw=FpgaBlue] (1.2,0.2) rectangle (3.4,0.8) node[midway] {$t_{\text{logic}}$};
  \fill[orange!20, draw=FpgaOrange] (3.4,0.2) rectangle (6.6,0.8) node[midway] {$t_{\text{wire}}$};
  \fill[FpgaLight, draw=FpgaBlue] (6.6,0.2) rectangle (7.5,0.8) node[midway] {$t_{su}$};
  \fill[green!20, draw=FpgaGreen] (7.5,0.2) rectangle (8,0.8);
  \draw[<->] (0,1.6) -- node[above] {period $T$} (8,1.6);
  \node[text=FpgaGreen!50!black] at (7.75,-0.6) {slack};
  \draw[->, FpgaGreen!50!black] (7.75,-0.4) -- (7.75,0.15);
\end{tikzpicture}
\caption{A path between two flip-flops. If the sum of the delays fits within the period, a positive slack remains}
\end{figure}

\subsection{How to Read the Report}

Open the Chip tab (the Real board group, which is collapsed by default), enter
the required value in the Frequency, MHz field, and click Place and build.
After the build, the Placement and routing panel shows the nextpnr result: the
maximum achievable frequency for each clock signal and its comparison with the
required one, while the Placement map on the chip shows where the cells were
placed. In Fig.~\ref{fig:n-chip}, the \code{clk} signal (27~MHz) has an
$f_{max}$ of about 339~MHz, which means a margin of roughly 12 times.

\shot{app-chip}{The nextpnr report in Logra (Running light project): the achievable frequency is 338.87~MHz against the required 27~MHz}{fig:n-chip}

If $f_{max}$ is lower than the required frequency, the circuit on the board
will work \emph{incorrectly, but without any error message}: the flip-flops
will capture values that have not had time to settle yet. That is why you must
look at the report after every build. A shortage of time can be fixed in
several ways: split long logic with a pipeline (Chapter~\ref{ch:sync}),
simplify the expression, reduce the number of inputs, or allow the synthesizer
to duplicate cells.

\begin{infobox}[The hold check]
Besides the setup slack, there is also a hold slack: the new value must not
arrive \emph{too early} after the edge, while the flip-flop is still holding the
old one. In an FPGA this margin is provided by the routing itself and by the
shared clock network, so engineers usually worry about it less.
\end{infobox}

\begin{diybox}[Find the critical path]
\textbf{Task.} Measure the maximum frequency of your project and see what happens when the required frequency is higher than the maximum.

\textbf{In Logra.} Open the Running light project and click Synthesize (Ctrl+Enter). On the Chip tab (Real board group), set the frequency to 27~MHz, click Place and build, and write down $f_{max}$. Then set a frequency above the recorded $f_{max}$ and build again.

\textbf{Expected result.} At 27~MHz, $f_{max}$ is about 339~MHz (a margin of 12 times); at a frequency above $f_{max}$, nextpnr reports that the requirement is not met, and the slack is negative.

\textbf{How to check.} The Placement and routing panel shows $f_{max}$ and a flag saying whether the requirement is met. Compare $T=1/f_{max}$ with the sum $t_{cq}+t_{\text{logic}}+t_{\text{wire}}+t_{su}$.

\stars{1} repeat: write down $f_{max}$ for the Running light. \stars{2} modify: lengthen the logic (for example, compare the counter with a 32-bit constant) and find out how $f_{max}$ changed. \stars{3} project: Security Panel stage---timing check; acceptance criterion: the panel build at 27~MHz meets timing with a margin of at least two times.
\end{diybox}

\section{Do Not Divide the Frequency---Use an Enable}

Beginners often ``divide the frequency'' like this: a flip-flop toggles on
every clock cycle, and its output is used as a new clock signal for another
part of the circuit. This is \emph{bad} practice. Such a ``clock'' does not
travel along the dedicated clock network but along ordinary wires, so it
reaches the flip-flops with different delays, and timing analysis becomes
unreliable.

The right way is to keep a single clock signal for the whole circuit and to
mark slow events with a \emph{clock enable pulse}. The divider counts clock
cycles and once every $N$ cycles produces a short \code{tick} pulse, and only
then do the flip-flops in other parts of the circuit change their values.

\begin{vcode}
reg [24:0] cnt = 0;
reg        tick = 0;                      // one cycle once per second
always @(posedge clk) begin
    if (cnt == 27_000_000 - 1) begin
        cnt  <= 0;
        tick <= 1;
    end else begin
        cnt  <= cnt + 1;
        tick <= 0;
    end
end

always @(posedge clk)                      // the same clk!
    if (tick) led <= ~led;                 // changes once per second
\end{vcode}

You used this technique in the projects of Chapter~\ref{ch:examples}. The
\code{tick} signal lasts exactly one clock cycle, but it controls any number
of flip-flops, and all of them stay synchronous with the single \code{clk}.

\begin{diybox}[One clock, many paces]
\textbf{Task.} Replace the ``divider-as-clock'' with an enable pulse: all flip-flops stay on one \code{clk}, and the slow events are marked by \code{tick}.

\textbf{In Logra.} Code tab: enter the fragment above (for the simulation, replace $27\,000\,000$ with 20 cycles). Waveforms: the signals \code{cnt}, \code{tick}, \code{led}.

\textbf{Expected result.} \code{tick} is 1 for exactly one cycle out of 20; \code{led} toggles only in the cycle with \code{tick = 1}.

\textbf{How to check.} On the waveform, the width of the \code{tick} pulse equals one period of \code{clk}. On the RTL tab, the same \code{clk} goes to all flip-flops, and there is no separate ``slow'' clock net.

\stars{1} repeat: the period of \code{tick} is 20 cycles. \stars{2} find the bug: take a counter bit as the clock for a second block and explain why timing analysis becomes unreliable because of this. \stars{3} project: Security Panel stage---paces; acceptance criterion: the siren beep and the exit delay both run from a single \code{clk} and the \code{tick} signals, and there is no second clock in the project.
\end{diybox}

\section{PLL: Multiplying the Frequency}

The Tang Nano 20K board has a 27~MHz crystal oscillator, but that is not a
limit. Inside the chip there is a \term{PLL} (phase-locked loop): a block that
makes another frequency from the input frequency $f_{\text{in}}$ according to
the formula

\[
f_{\text{out}} = f_{\text{in}}\cdot\frac{M}{N},
\]

where $M$ and $N$ are integer multiplication and division factors (the exact
allowed ranges are in the chip documentation). In this way, from 27~MHz you
can get, for example, 108~MHz for fast logic or 74.25~MHz for an HDMI video
signal. A PLL also cleans the clock signal of jitter and can produce several
phases.

A PLL is connected as a ready-made primitive (in Gowin it is called
\code{rPLL}), and its output goes to the global clock network. After that, the
circuit behind the PLL counts as a new \term{clock domain} with its own
frequency.

\begin{diybox}[Choose a PLL on paper and check the timing]
\textbf{Task.} Choose $M$ and $N$ so that 27~MHz becomes 108~MHz and 74.25~MHz, and estimate whether your logic can handle such a frequency.

\textbf{In Logra.} Do the calculations with the formula $f_{\text{out}}=f_{\text{in}}M/N$. Then, on the Chip tab, enter 108 in the Frequency, MHz field and rebuild your project.

\textbf{Expected result.} For 108~MHz, $M/N=4$ works; for 74.25~MHz, $M/N=2.75=11/4$. If the project's $f_{max}$ is below 108~MHz, nextpnr reports that the requirement is not met, and the logic will have to be made faster.

\textbf{How to check.} Substitute the result into the formula; in the Placement and routing panel, compare $f_{max}$ with the required frequency.

\stars{1} repeat: 108~MHz. \stars{2} modify: choose $M$ and $N$ for 148.5~MHz (an HDMI 1080p pixel clock; $M/N=5.5$). \stars{3} project: Security Panel stage---clock frequency; acceptance criterion: $M$, $N$, and the expected frequency are written down, and the panel's $f_{max}$ is higher than this frequency.
\end{diybox}

\section{Multiple Clock Domains}

When a circuit has two different clock signals, a new problem appears: a signal
that comes from one domain into another may change right \emph{during} the
edge of the receiving clock. Then the setup/hold condition is violated, and
the flip-flop enters a metastable state (Chapter~\ref{ch:ff}): its output hangs
between 0 and 1 and decides on its own after an undefined time. If the
unlucky value gets straight into the logic, the circuit breaks.

\subsection{The Two-Flip-Flop Synchronizer}

A simple solution is to pass the signal through a chain of \emph{two}
flip-flops in the receiving domain. The first one may become metastable, but
it has a whole clock cycle to ``decide,'' and the second one receives an
already stable value. The probability that metastability survives a whole
clock cycle is negligibly small.

\vfile{\codepath/10_clocks/sync2ff.v}{code/10\_clocks/sync2ff.v}

\shot{n-sync-rtl}{The synchronizer in Logra: two flip-flops one after another, followed by the debounce filter}{fig:n-sync-rtl}

The same synchronizer is needed for any asynchronous signal: a button, or a
signal from another chip (for example, the input of a serial-interface
receiver, Chapter~\ref{ch:uart}).

\begin{warnbox}
A two-flip-flop synchronizer works only for \emph{a single bit}. If you pass a
multi-bit number (for example, a counter), the bits may arrive with different
delays, and the receiving side will briefly see a value that never existed.
Buses call for other techniques: pass the data one bit at a time with a
handshake (request---acknowledge), use a \emph{Gray code}, where only one bit
changes at any step (Chapter~\ref{ch:codes}), or use dual-clock FIFOs whose
pointers are sent across in Gray code.
\end{warnbox}

\subsection{Contact Bounce}

A mechanical button does not close just once when pressed: the contacts bounce
for several milliseconds, and the signal chatters between 0 and 1. A digital
circuit perceives this as many presses. The fix is a counter that accepts a
new value only if it holds for $N$ clock cycles in a row.

\shot{n-debounce-wave}{The debounce filter: the input \code{in} jumps 1-0-1-0-1, but the output \code{out} switches only after the value has held for several clock cycles}{fig:n-debounce}

The waveform (Fig.~\ref{fig:n-debounce}) shows two important things. The signal
\code{s} after the synchronizer repeats \code{in} with a delay of two clock
cycles. And the output \code{out} waits until \code{s} settles: it ignores
false short pulses.

\begin{diybox}[Synchronizer and bounce]
\textbf{Task.} Pass an asynchronous button through a two-flip-flop synchronizer and a debounce filter, and make sure the output changes only once on a burst of noise.

\textbf{In Logra.} Code tab: \code{sync2ff.v} from the book; Synthesize; the RTL tab.

\textbf{Expected result.} Two flip-flops go one after another, followed by the debounce filter; the output \code{out} changes only after the value has held for several clock cycles.

\textbf{How to check.} Waveforms: apply the series 1-0-1-0-1 to \code{in} and then a long 1; the signal \code{s} repeats \code{in} with a two-cycle delay, and \code{out} ignores the short pulses.

\stars{1} repeat: the series 1-0-1-0-1. \stars{2} modify: double the number of filter cycles and check that short pulses are still rejected. \stars{3} project: Security Panel stage---code entry; acceptance criterion: on the waveform, three presses of the code button with bounce on the input give exactly three \code{code\_step} pulses.
\end{diybox}

\section{Reset}

The reset signal must also be synchronous. An asynchronous reset acts
instantly, but its \emph{release} at a random moment can violate the setup time
of some of the flip-flops and bring the circuit into different states. The
right approach is to apply the reset to the flip-flops synchronously, and to
pass an external reset button through the same two-flip-flop synchronizer. The
examples in this book use a synchronous reset everywhere: it is simpler and
more reliable, and a chip of this type does not need separate circuits for it.

\begin{diybox}[Synchronous reset of the panel]
\textbf{Task.} Apply the reset to the panel synchronously: the external button passes through a synchronizer, and all registers are reset on the clock edge.

\textbf{In Logra.} Code tab: first the \code{guard\_panel} module from Chapter~\ref{ch:bridge} (it has \code{if (rst)} inside \code{always @(posedge clk)}), then put two flip-flops in front of \code{rst}.

\textbf{Expected result.} The reset acts on the nearest clock edge, not instantly; all the panel's flip-flops are reset in the same cycle.

\textbf{How to check.} Waveforms: raise \code{rst} between edges and see when \code{state} returns to OFF. RTL: there are no elements with an asynchronous reset, only the synchronous \code{SRST}.

\stars{1} repeat: find the reset cycle on the waveform. \stars{2} find the bug: describe the reset asynchronously (\code{always @(posedge clk or posedge rst)}) and explain why releasing the reset is dangerous. \stars{3} project: Security Panel stage---reset; acceptance criterion: on reset, in any state, the panel ends up in OFF within one clock cycle after the synchronizer.
\end{diybox}

\section{Questions and Problems}
\label{sec:tasks-clocks}

The answers are collected in Appendix~\ref{app:answers}.

\begin{taskbox}[Chapter problems]
\begin{enumerate}
  \item \stars{1}~A path between flip-flops: $t_{cq}=0.5$~ns, logic $4$~ns, wires $3$~ns,
        setup time $0.5$~ns. What must the period be? What is $f_{max}$?
  \item \stars{1}~The clock period is 10~ns, and the total path delay is 8.2~ns. What is the slack?
        What does it become at 150~MHz?
  \item \stars{2}~Describe a divider that produces a \code{tick} pulse at 1~kHz
        from 27~MHz. What value must the counter count up to?
  \item \stars{1}~A PLL receives 27~MHz, $M=4$, $N=1$. What is the output frequency?
  \item \stars{2}~Why can't you pass an 8-bit counter from one domain to another
        through eight separate synchronizers?
  \item \stars{2}~A debounce filter must reject noise of up to 5~ms at a clock frequency of
        27~MHz. What is the value of $N$?
  \item \stars{2}~Find the bug. According to the report, the longest path between flip-flops takes 12~ns.
        The designer reasons: ``I'll set the PLL to 100~MHz, since the chip's frequency allows it.''
        What will actually happen? What is the slack?
  \item \stars{2}~Find the bug. In a circuit, a slow block is clocked by the output of a flip-flop
        ``divide by 2'' (\code{always @(posedge slow\_clk)}). In simulation everything works, but on the board
        there are occasional failures. Name the cause and rewrite the circuit with an enable pulse.
  \item \stars{3}~Security Panel project, stage ``timing and reset.'' Build the panel from Chapter~\ref{ch:bridge},
        feed the buttons through a synchronizer and debouncing, and use a synchronous reset. Acceptance
        criterion: on the Chip tab, $f_{max}$ is at least 54~MHz (twice 27~MHz), and on the waveform
        a burst of bounce on the button leads to exactly one transition.
  \item \stars{3}~Project. Build a counter with a clock of $f=27$~MHz that toggles an LED every 0.5~s,
        using only an enable pulse (no second clock). Acceptance criterion:
        in the RTL there are no nets where a flip-flop output goes to a clock input, and $f_{max}$ in the report
        is at least 100~MHz.
\end{enumerate}
\end{taskbox}
